\chapter{Вывод проекторной фазы и точная граница доказанного}

Проекторное быстрое переключение имеет физический смысл только в том случае, если
$\mathbb{RP}^2$ возникает как пространство минимумов родительского
функционала. Поэтому необходимо отделить топологию уже выбранного
проектора от динамического вывода самой упорядоченной фазы.

\section{Обратный аудит проектных оснований}

В Томе V семейный проектор записан как
\begin{equation}
 P=nn^T,\qquad n\in S^2,\qquad n\sim-n.
 \label{eq:v6-rp2-orbit}
\end{equation}
Отсюда строго следует, что пространство проекторов ранга один равно
$\mathbb{RP}^2$. Для этой орбиты доказаны касательный модуль, проекторный
ёж, спинорно-хопфов мост и локальные классы $+15/-15$.

Но проверки о касательном модуле и угле состояния явно начинались с уже
выбранного $P$ или $\rho$. Действие Ходжа на $M_{300}(\mathbb C)$ имеет
минимум полного ранга $X=I_3$ и оставляет семейный проектор вне полного
функционала. Положительный квадрат центрированной связности выбирает
нулевую связность. Аффинный проектор $P_1$ является фиксированной
синглетной линией четырёхмерного меню, а не подвижной семейной орбитой.

Существующее модулярное состояние также не выводит нужную фазу. Его
минимальная высота имеет кратность шесть, поэтому при
$\beta\to\infty$ получается нормированный проектор ранга шесть; внутри
каждого семейного триплета состояние остаётся изотропным.

\begin{theorem}[Статус $\mathbb{RP}^2$ в текущем родителе]
Проект строго вывел $\mathbb{RP}^2$ как пространство конфигураций
условно выбранной семейной проекторной фазы, но ни один существующий
родительский функционал не имеет доказанного множества минимумов,
равного этой орбите.
\end{theorem}

\section{Признанный вариационный контроль}

В теории Ландау--де Жена вещественный симметричный бесследовый тензор
$Q$ имеет стандартную объёмную плотность
\begin{equation}
 f_B(Q)=\frac a2\Tr Q^2-\frac b3\Tr Q^3
       +\frac c4(\Tr Q^2)^2,qquad c>0.
 \label{eq:v6-ldg-control}
\end{equation}
В упорядоченной области минимумы имеют вид
\begin{equation}
 Q=s_\ast\left(nn^T-\frac13I_3\right),\qquad n\sim-n,
 \label{eq:v6-ldg-rp2-minima}
\end{equation}
и образуют $\mathbb{RP}^2$. Это признанный пример вариационного рождения
проективной орбиты; см. \path{arXiv:0812.3131},
\path{arXiv:0808.1870} и \path{arXiv:2010.07786}.

При $a=-1$, $b=c=1$ получаем $s_\ast=3/2$. Машинный гессиан в
пятимерном пространстве бесследовых симметричных матриц имеет две нулевые
касательные и три положительные поперечные моды. Но коэффициенты $a,b,c$
не фиксируются одной симметрией. Поэтому этот потенциал является
контрольным механизмом, а не готовым родительским выводом проекта.

\section{Носитель уже присутствует}

Вещественный семейный фактор естественно содержит пространство
\begin{equation}
 \mathcal D_{\mathbb R,3}
 =\{R=R^T\geq0:\Tr R=1\}.
 \label{eq:v6-real-state-space}
\end{equation}
Его центр $I_3/3$ изотропен, а крайние точки $P=nn^T$ образуют
$\mathbb{RP}^2$. Следовательно, новую трёхмерную кинематику вводить не
нужно: недостаёт вариационного механизма, который нарушает трёхкратное
вырождение состояния.

Важно не потребовать лишнего. Уже семейство смешанных состояний
\begin{equation}
 R_x=xP+\frac{1-x}{2}(I_3-P),\qquad x\ne\frac13,
 \label{eq:v6-uniaxial-mixed-state}
\end{equation}
имеет одно простое и одно двукратное собственное значение. Его
ориентационная орбита также равна $SO(3)/O(2)=\mathbb{RP}^2$. Значит,
для рождения проекторной топологии достаточно спектрального расщепления
$3\to1+2$; полная очистка $R^2=R$ не обязательна.

\begin{proposition}[Уточнённый недостающий мост]
Проекту недостаёт не пространства $\mathbb{RP}^2$, а родительского
члена, который переводит изотропное состояние $I_3/3$ в одноосную
страту \eqref{eq:v6-uniaxial-mixed-state}, и доказательства, что это
семейное состояние является параметром порядка возникающей геометрии.
\end{proposition}

\begin{nogocriterion}[Запрет подмены геометрической фазы]
Нельзя считать $\mathbb{RP}^2$ выведенным из родителя, пока не получены
общий функционал на $\mathcal D_{\mathbb R,3}$, потеря глобальной
устойчивости изотропного состояния, одноосные минимумы и отображение
семейного параметра порядка в геометрические наблюдаемые.
\end{nogocriterion}

\begin{protocol}
\begin{enumerate}
 \item $\mathbb{RP}^2$ как орбита выбранного проектора:
 \textbf{строго доказано};
 \item проекторный ёж и классы $+15/-15$:
 \textbf{строги условно на фазе};
 \item $\mathbb{RP}^2$ как минимумы текущего действия:
 \textbf{не выведено};
 \item признанный вариационный аналог:
 \textbf{есть};
 \item новый трёхмерный носитель:
 \textbf{не требуется};
 \item полная очистка до ранга один:
 \textbf{достаточна, но не необходима};
 \item минимальное условие:
 \textbf{расщепление спектра $3\to1+2$};
 \item геометрическая атрибуция состояния:
 \textbf{открыта};
 \item следующая проверка:
 \textbf{энтропийно-ориентационный переход вещественного состояния}.
\end{enumerate}
\end{protocol}

Машинный сертификат сохранён в
\path{s2t_v6_rp2_geometric_phase_derivation_gate_results.json}.